\section{Exact validation of rational descriptions}
The rational codes below require an exact positivity test. We record
the arithmetic fact in the form used by both encoders.

\begin{lemma}[Reduced Schur entries and bit length]\label{lem:schurminors67}
Let $A$ be a Hermitian matrix of order $q$ with Gaussian-integer entries,
whose real and imaginary parts have absolute value at most $2^h$.
In exact Schur elimination, let $I$ be the ordered indices of the
positive pivots already eliminated. Every remaining entry is
\begin{equation}\label{eq:borderminor67}
 S_{ij}=\frac{\det A[I\mathbin{\|}i,I\mathbin{\|}j]}{\det A[I,I]},
\end{equation}
where $I\mathbin{\|}i$ appends $i$ in that order, and the empty
principal determinant is one. Reduced numerators and denominators
have $O(q(h+\log(q+1)))$ bits. Hermitian positive semidefiniteness is
therefore decidable by this elimination in polynomial bit time and
space in the explicit input length. No optimized polynomial exponent
is asserted.
\end{lemma}
\begin{proof}
When the leading pivot $t$ is positive, block congruence reduces
\[
 \begin{pmatrix}t&b^*\\ b&C\end{pmatrix}\ge0
 \quad\Longleftrightarrow\quad C-bb^*/t\ge0.
\]
A negative pivot rejects. At a zero diagonal, each $2\times2$ principal
minor forces the corresponding row and column to vanish; otherwise
reject, and if they vanish delete that row and column. Positive pivots
make $A[I,I]$ positive definite. The bordered determinant identity
\[
 \det\begin{pmatrix}A[I,I]&A[I,j]\\A[i,I]&A_{ij}\end{pmatrix}
 =\det A[I,I]\,(A_{ij}-A[i,I]A[I,I]^{-1}A[I,j])
\]
proves \eqref{eq:borderminor67}, including after zero rows are skipped.
The specified ordering prevents a spurious permutation sign.

A minor of order $k\le q$ has determinant modulus at most
$k^{k/2}(\sqrt2\,2^h)^k$, by the product of its row norms.
Its real and imaginary parts are integers with
$O(k(h+\log(k+1)))$ bits. The principal denominator is a positive
integer. Cancelling common factors cannot enlarge these lengths.
Each arithmetic operation on reduced Gaussian-rational entries thus
has polynomial operand size. There are polynomially many operations
and at most $q^2$ stored entries; ordinary multiplication, division and
the Euclidean algorithm give the stated bit bounds. For rational
rather than integral input, first clear a common positive denominator.
Its bit length is at most the sum of the encoded denominator lengths,
so this preliminary step also has polynomial cost.
\end{proof}
